\chapter{Phương pháp nghiên cứu}
\label{chap:methodology}

\section{Thiết kế nghiên cứu}

Nghiên cứu sử dụng thiết kế thực nghiệm ghép cặp trên một kho mẫu dùng chung. Với mỗi
bài toán, mỗi token cap và mỗi chỉ số mẫu, hệ thống sinh đúng một quỹ đạo cho lần lặp
chính. Các cấu hình cố định và chính sách thích ứng sau đó được phát lại trên các tiền
tố của cùng kho mẫu. Thiết kế này loại bỏ khác biệt do các policy tình cờ nhận các lần
lấy mẫu khác nhau, đồng thời giảm số lần suy luận GPU cần thiết.

Ba lớp bằng chứng được tách biệt. Phân tích Qwen3-1.7B/GSM8K là bằng chứng lịch sử dùng
để hình thành giả thuyết. Train và Validation của Qwen3-0.6B dùng để xây dựng ngưỡng,
chọn policy và quyết định pivot. Test-100 của GSM8K và MATH-500 là đánh giá giữ lại
cho các claim cuối. Không điều chỉnh policy sau khi xem kết quả Test.

\section{Ký hiệu và không gian hành động}

Gọi $x_i$ là đề bài và $y_i$ là đáp án chuẩn của bài $i$. Một profile cố định là

\begin{equation}
  p=(k,L),
\end{equation}

trong đó $k$ là số quỹ đạo độc lập và $L$ là giới hạn token sinh trên mỗi quỹ đạo.
Không gian breadth chung là $k\in\{1,2,4\}$. Với GSM8K, không gian depth là
$L\in\{512,1024,2048\}$; với MATH-500, không gian depth là
$L\in\{1024,2048,4096\}$. Mỗi dataset vì vậy có chín profile cố định.

Ngân sách danh nghĩa của profile là

\begin{equation}
  B(p)=kL.
\end{equation}

Nếu $t_{ij}(L)$ là số token thật sự sinh cho mẫu $j$ của bài $i$, chi phí thực tế là

\begin{equation}
  C_i(p)=\sum_{j=1}^{k}t_{ij}(L), \qquad 0\leq t_{ij}(L)\leq L.
  \label{eq:realized-cost}
\end{equation}

Mọi quỹ đạo policy đã yêu cầu đều xuất hiện trong tổng ở Phương trình
\ref{eq:realized-cost}, kể cả quỹ đạo không đóng góp phiếu thắng. Latency được ghi lại
như một metric phụ, nhưng không dùng làm cost chính vì phụ thuộc mạnh vào phần cứng và
trạng thái hệ thống.

\section{Sinh mẫu ghép cặp và tái sử dụng tiền tố}

Với mỗi cặp $(i,L)$, hệ thống sinh bốn mẫu có chỉ số $0,1,2,3$. Profile $k=1$ dùng mẫu
0; $k=2$ dùng mẫu 0--1; $k=4$ dùng mẫu 0--3. Do đó, một bài chỉ cần 12 lần sinh cho ba
giá trị $L$, thay vì sinh lại độc lập cho từng cặp $(k,L)$.

Seed của một mẫu được xác định bởi

\begin{equation}
  s_{ijr}=\operatorname{first63bits}\!\left(
  \operatorname{SHA256}(s_0\,|\,i\,|\,r\,|\,j)\right),
\end{equation}

với $s_0$ là global seed, $r$ là repetition ID và $j$ là sample ID. Seed cố ý không
phụ thuộc vào $k$ hoặc $L$. Cùng sample ID vì vậy dùng nguồn ngẫu nhiên ghép cặp giữa
các token cap, còn việc không đưa $k$ vào seed bảo đảm các profile breadth là tiền tố
lồng nhau. Global seed của thực nghiệm chính là 42 và chỉ có một repetition. Vì vậy,
khoảng bất định thống kê phản ánh biến thiên giữa các bài, không phản ánh biến thiên
giữa nhiều lần sinh độc lập.

\section{Parse, chuẩn hóa và tổng hợp đáp án}

Parser chỉ chấp nhận đáp án cuối trong các định dạng đã đăng ký: trường final answer,
marker của GSM8K, biểu thức boxed hoặc định dạng số rõ ràng tùy dataset. Hệ thống không
dùng quy tắc ``lấy số cuối cùng'' từ toàn bộ reasoning. Mỗi mẫu được phân loại thành
một trong bốn trạng thái: bị cắt trong thinking, bị cắt sau thinking, hoàn tất nhưng
không parse được, hoặc hoàn tất với đáp án hợp lệ.

Gọi $V_i(p)$ là dãy đáp án chuẩn hóa hợp lệ trong $k$ mẫu. Dự đoán profile là mode của
$V_i(p)$. Nếu có hòa phiếu, mode xuất hiện ở sample ID thấp nhất được chọn. Mẫu invalid
không bỏ phiếu; nếu $V_i(p)=\varnothing$, profile được tính sai. Độ đồng thuận trên đáp
án hợp lệ là

\begin{equation}
  G_i(p)=\frac{\max_a \sum_{v\in V_i(p)}\mathbf{1}[v=a]}{|V_i(p)|},
\end{equation}

khi $|V_i(p)|>0$. Tỷ lệ đáp án hợp lệ được tính trên toàn bộ $k$ mẫu. Correctness
$A_i(p)\in\{0,1\}$ là so khớp chính xác sau chuẩn hóa giữa dự đoán tổng hợp và $y_i$.

GSM8K sử dụng chuẩn hóa số; MATH-500 sử dụng parser exact cho biểu thức LaTeX thường
gặp. Trong pilot lịch sử 1.7B, mục GSM8K có ID \code{gsm8k-test-01309} được giữ
trong raw artifact để bảo toàn provenance nhưng bị loại khỏi metric do cờ chất
lượng dữ liệu. Mục này không thuộc manifest Test-100 chính của thực nghiệm 0.6B.
Nhãn gốc không bị sửa nhằm tránh thay đổi benchmark theo kết quả mô hình.

\section{Các họ policy so sánh}

\subsection{Profile cố định}

Policy cố định áp dụng cùng một $(k,L)$ cho mọi bài. Chín profile tạo các anchor để
kiểm tra tác động riêng của breadth và token cap. Chúng cũng là chuẩn cần thiết để
xác định liệu một policy thích ứng có thực sự tạo điểm accuracy--token tốt hơn hay chỉ
chọn lại một profile cố định tốn kém.

\subsection{Depth-only}

Depth-only luôn dùng $k=1$. Số prompt token của các bài Train được dùng để tính hai
ngưỡng tertile bao hàm. Bài ở tertile thấp, giữa và cao lần lượt nhận ba giá trị $L$
đã sắp tăng dần. Quy tắc này không nhìn output và không ước lượng xác suất đúng. Nó là
baseline thích ứng theo đầu vào, dùng để kiểm tra giá trị tăng thêm của selective
breadth.

\subsection{Breadth-only agreement stopping}

Breadth-only giữ một $L$ cố định. Policy luôn mua hai mẫu đầu. Nếu cả hai parse hợp lệ
và có đáp án chuẩn hóa giống nhau, policy dừng ở $k=2$; nếu không, policy mua tiếp mẫu
2 và 3 để đạt $k=4$. Không có nhánh dừng ở $k=1$. Tất cả mẫu đã mua được tính vào chi
phí và dự đoán cuối được tổng hợp bằng cùng quy tắc valid-only.

\subsection{Heuristic joint depth-first}

Heuristic joint trước hết dùng prompt-token tertile từ Train để chọn $L$. Tại $L$ đã
chọn, policy đọc sample pool theo thứ tự. Policy dừng ở $k=1$ khi mẫu đầu có đáp án
cuối hợp lệ và không bị cắt; nếu điều kiện này không đạt, policy kiểm tra hai mẫu đầu
và dừng ở $k=2$ khi cả hai hợp lệ, giống nhau; mọi trường hợp còn lại lấy đủ $k=4$.

\begin{table}[htbp]
  \centering
  \caption{Thông tin được phép dùng tại thời điểm ra quyết định.}
  \label{tab:policy-signals}
  \begin{tabular}{p{0.52\textwidth}p{0.16\textwidth}p{0.16\textwidth}}
    \toprule
    Tín hiệu & Được dùng & Bị cấm \\
    \midrule
    Số token của prompt hiện tại & Có & -- \\
    Đáp án parse được từ mẫu đã sinh & Có & -- \\
    Trạng thái truncated của mẫu đã sinh & Có & -- \\
    Đáp án chuẩn của bài hiện tại & -- & Có \\
    Nhãn/profile do oracle chọn & -- & Có \\
    Mẫu tương lai hoặc nhánh phản thực tế & -- & Có \\
    Kết quả Test để chỉnh ngưỡng & -- & Có \\
    \bottomrule
  \end{tabular}
\end{table}

Joint policy này không phải controller học máy. Ngưỡng tối thiểu để cho phép huấn
luyện một stage classifier là 30 sự kiện dương trên Train; dữ liệu không đạt gate nên
nhánh learned controller bị dừng. Quyết định này giảm nguy cơ overfit một tập sự kiện
quá nhỏ.

\section{Oracle và giới hạn diễn giải}

Oracle được tính trên sample pool đã sinh để đo headroom, không được triển khai như
một policy. Với

\begin{equation}
  \mathcal{S}_i=\{p:A_i(p)=1\},
\end{equation}

oracle chọn profile đúng có thứ tự từ điển nhỏ nhất theo

\begin{equation}
  r_i(p)=\big(C_i(p),T_i(p),B(p),k,L\big),
\end{equation}

trong đó $T_i(p)$ là latency. Nếu $\mathcal{S}_i$ rỗng, oracle chọn profile quan sát
có thứ tự chi phí nhỏ nhất thay vì loại bài khỏi mẫu số. Quy tắc này định nghĩa rõ chi
phí cho bài không giải được.

Oracle nhìn đáp án chuẩn và chọn sau khi đã thấy mọi cấu hình, nên hưởng lợi từ các
mẫu đúng do may mắn trong một lần sinh. Vì vậy, oracle chỉ là cận trên cơ hội của pool
một seed. Khoảng cách giữa policy và oracle không phải bằng chứng rằng một controller
thực tế có thể đạt mức oracle, cũng không phải ước lượng không chệch của compute tối
ưu trong phân phối.

\section{Protocol chọn và đóng băng policy}

Protocol gồm bốn bước. Một là tạo split manifest tất định và kiểm tra hash câu hỏi,
đáp án. Hai là dùng Train để tính tertile, tạo event tuần tự và kiểm tra gate của
learned controller. Ba là dùng Validation để chọn điểm vận hành trong các policy đã
định nghĩa trước. Bốn là tuần tự hóa policy ID, ngưỡng, action space, prompt, parser,
seed và software hash trước khi chạy Test.

Trên Validation GSM8K, breadth-only tại $L=1024$ và joint heuristic cùng đạt 15/20,
nhưng breadth-only dùng trung bình 2095,85 token so với 2149,05; breadth-only vì vậy
được chọn. Trên Validation MATH-500, các policy thích ứng bị các profile cố định chi
phối; profile $k=2,L=4096$ đạt 13/20 với 4700,50 token, trong khi joint đạt 10/20 với
5935,55 token. Profile cố định được chọn cho điểm vận hành chính của MATH-500. Đây là
bằng chứng chọn mô hình trên Validation, không phải kết quả Test.

Tổng cộng 14 policy/comparator đã đăng ký được phát lại trên mỗi Test-100. Test không
dùng để đổi policy chính, ngưỡng hoặc parser. Pareto frontier trên Test chỉ mang tính
mô tả cho tập policy đã đóng băng.

\section{Metric và phân tích thống kê}

Hai metric chính là exact-answer accuracy và số token sinh thực tế trung bình. Các
metric phụ gồm $kL$, latency, tỷ lệ truncated, tỷ lệ đáp án hợp lệ và token trên mỗi
bài đúng. Khoảng tin cậy 95\% cho accuracy dùng Wilson score interval, phù hợp hơn xấp
xỉ chuẩn khi mẫu nhỏ hoặc accuracy gần biên.

So sánh accuracy giữa hai policy trên cùng bài dùng exact McNemar test. Với một họ so
sánh đã lên kế hoạch, $p$-value được hiệu chỉnh Holm. Chênh lệch accuracy và token
được bổ sung bằng bootstrap ghép cặp ở cấp bài, tức mỗi lần bootstrap lấy lại các bài
và giữ cặp outcome giữa hai policy. Không bootstrap từng mẫu sinh vì đơn vị suy luận
và đơn vị chia split là bài toán.

Một policy $a$ chi phối policy $b$ nếu accuracy của $a$ không thấp hơn, token trung
bình không cao hơn và có ít nhất một bất đẳng thức nghiêm. Biên Pareto gồm các policy
không bị chi phối. Quan hệ này mô tả hai metric quan sát; nó không tự cung cấp kiểm
định ý nghĩa thống kê và không thay thế bảng so sánh ghép cặp.

\section{Đe dọa hiệu lực được xử lý trong thiết kế}

Rò rỉ Test được kiểm soát bằng manifest và policy freeze. Sai lệch do sample khác nhau
được giảm bằng nested shared pool. Sai lệch parser được giảm bằng parser versioned,
không last-number fallback, và giữ raw response. Sai lệch chi phí được giảm bằng actual
generated tokens và việc tính phí mọi mẫu đã mua. Tuy vậy, thiết kế vẫn có các giới
hạn không thể loại bỏ trong phạm vi đề tài: một repetition, một mô hình chính, hai bộ
dữ liệu toán, và không có đối chứng decoding non-thinking/budget-forcing. Các giới hạn
này phải được giữ khi diễn giải kết quả ở các chương sau.
